\documentclass[11pt]{article}

\usepackage[margin=1in]{geometry}
\usepackage{amsmath}
\usepackage{amssymb}
\usepackage{amsthm}

\DeclareMathOperator{\rank}{rank}

\theoremstyle{definition}
\newtheorem{problem}{Problem}
\newtheorem*{remark}{Remark}
\newtheorem*{corollary}{Corollary}

\title{IE 5311-001 (D01) Principles of Optimization\\[2pt]
       \large Class 01, Slide 7 Self-Check Problems}
\author{Lidivine Kengne}
\date{\today}

\begin{document}

\maketitle

%======================================================================
\section{Coordinates Under Two Bases}
%======================================================================

\begin{problem}
Consider the vector space of degree-2 polynomials with real coefficients,
and the two bases
\[
  \mathcal{B}_1 = \{1,\; x,\; x^2\},
  \qquad
  \mathcal{B}_2 = \{1+x,\; x+x^2,\; 1+x^2\}.
\]
For the polynomial $p(x) = 3 + 5x + 4x^2$, find its coefficients under
$\mathcal{B}_1$ and under $\mathcal{B}_2$, respectively.
\end{problem}

\subsection*{Setup}

Write $\mathcal{P}_2$ for the space of real polynomials of degree at most
two. Addition and scalar multiplication act coefficient by coefficient,
so $\dim \mathcal{P}_2 = 3$.

A vector and its coordinate list are different objects. The polynomial
$p$ is fixed. The list of numbers describing $p$ changes with the basis.
For a basis $\mathcal{B} = \{v_1, v_2, v_3\}$ we write
$[p]_{\mathcal{B}} = (a_1, a_2, a_3)^{\top}$ to mean
$p = a_1 v_1 + a_2 v_2 + a_3 v_3$.

\subsection*{Coordinates under $\mathcal{B}_1$}

The basis $\mathcal{B}_1 = \{1, x, x^2\}$ is the standard one, so the
coefficients are read directly from $p$:
\[
  p(x) = 3 \cdot 1 \;+\; 5 \cdot x \;+\; 4 \cdot x^2,
  \qquad\text{hence}\qquad
  [p]_{\mathcal{B}_1} = \begin{pmatrix} 3 \\ 5 \\ 4 \end{pmatrix}.
\]

\subsection*{Coordinates under $\mathcal{B}_2$}

Seek scalars $a, b, c \in \mathbb{R}$ such that
\[
  a(1+x) \;+\; b(x+x^2) \;+\; c(1+x^2) \;=\; 3 + 5x + 4x^2 .
\]
Expand the left-hand side and collect powers of $x$:
\[
  a(1+x) + b(x+x^2) + c(1+x^2)
  \;=\; \underbrace{(a+c)}_{\text{constant}}
  \;+\; \underbrace{(a+b)}_{x}\, x
  \;+\; \underbrace{(b+c)}_{x^2}\, x^2 .
\]
Two polynomials are equal if and only if their coefficients agree power
by power, which gives the linear system
\[
  \begin{cases}
    a + c = 3, \\
    a + b = 5, \\
    b + c = 4.
  \end{cases}
\]
From the first equation, $c = 3 - a$. Substituting into the third gives
$b = 4 - c = 1 + a$. Substituting that into the second gives
$a + (1+a) = 5$, hence $a = 2$, and back-substitution yields $b = 3$ and
$c = 1$. Therefore
\[
  [p]_{\mathcal{B}_2} = \begin{pmatrix} 2 \\ 3 \\ 1 \end{pmatrix}.
\]

\subsection*{Verification}

\[
  2(1+x) + 3(x+x^2) + 1(1+x^2)
  = (2+1) + (2+3)x + (3+1)x^2
  = 3 + 5x + 4x^2 = p(x). \checkmark
\]

\subsection*{Change-of-basis matrix}

Collect the $\mathcal{B}_2$ vectors, written in $\mathcal{B}_1$
coordinates, as the columns of a matrix:
\[
  M =
  \begin{pmatrix}
    1 & 0 & 1 \\
    1 & 1 & 0 \\
    0 & 1 & 1
  \end{pmatrix},
  \qquad
  \text{columns } \;1+x,\;\; x+x^2,\;\; 1+x^2 ,
\]
whose rows correspond to the coefficients of $1$, $x$ and $x^2$. Then
\[
  M \, [p]_{\mathcal{B}_2} = [p]_{\mathcal{B}_1},
  \qquad\text{equivalently}\qquad
  [p]_{\mathcal{B}_2} = M^{-1} [p]_{\mathcal{B}_1}.
\]
Cofactor expansion along the first row gives
\[
  \det M = 1\,(1\cdot 1 - 0\cdot 1) - 0\,(1\cdot 1 - 0\cdot 0)
         + 1\,(1\cdot 1 - 1\cdot 0) = 2 \neq 0 ,
\]
which confirms that $\mathcal{B}_2$ is a basis rather than an arbitrary
set of three polynomials. A direct check confirms the answer:
\[
  M \begin{pmatrix} 2 \\ 3 \\ 1 \end{pmatrix}
  = \begin{pmatrix} 2 + 0 + 1 \\ 2 + 3 + 0 \\ 0 + 3 + 1 \end{pmatrix}
  = \begin{pmatrix} 3 \\ 5 \\ 4 \end{pmatrix}
  = [p]_{\mathcal{B}_1}. \checkmark
\]

\begin{remark}
In linear programming a basic feasible solution is a choice of basis, and
the simplex method is a rule for moving from one basis to the next. The
determinant check performed here is the same nonsingularity condition
imposed on a simplex basis matrix.
\end{remark}

%======================================================================
\section{Kernel of a Transpose Product}
%======================================================================

\begin{problem}
Given any matrix $A$, prove that $\ker(A^{\top}\!A) = \ker(A)$.
\end{problem}

\subsection*{Setup}

Let $A \in \mathbb{R}^{m \times n}$. Recall
\[
  \ker(A) = \{\, x \in \mathbb{R}^{n} : Ax = 0 \,\}.
\]
Since $A^{\top}\!A \in \mathbb{R}^{n \times n}$, both kernels are
subspaces of $\mathbb{R}^{n}$, so the comparison is well posed. Proving
equality of two sets requires two inclusions.

\begin{proof}
\textbf{Step 1: $\ker(A) \subseteq \ker(A^{\top}\!A)$.}

Let $x \in \ker(A)$, so $Ax = 0$. Multiplying on the left by $A^{\top}$
and using associativity,
\[
  (A^{\top}\!A)\,x = A^{\top}(Ax) = A^{\top} 0 = 0 ,
\]
hence $x \in \ker(A^{\top}\!A)$.

\medskip
\textbf{Step 2: $\ker(A^{\top}\!A) \subseteq \ker(A)$.}

Let $x \in \ker(A^{\top}\!A)$, so $A^{\top}\!A x = 0$. Multiply on the
left by $x^{\top}$:
\[
  0 = x^{\top}\!\left(A^{\top}\!A x\right)
    = \left(x^{\top} A^{\top}\right)(A x)
    = (Ax)^{\top}(Ax)
    = \left\| Ax \right\|_{2}^{2} .
\]
Writing $y = Ax \in \mathbb{R}^{m}$, this says
$y_1^2 + y_2^2 + \cdots + y_m^2 = 0$. A sum of squares of real numbers
vanishes if and only if every term vanishes, so $y_i = 0$ for all $i$,
that is $Ax = 0$ and $x \in \ker(A)$.

\medskip
Both inclusions hold, therefore $\ker(A^{\top}\!A) = \ker(A)$.
\end{proof}

\begin{remark}[The real field is essential]
Step 2 uses the fact that a sum of squares of \emph{real} numbers
vanishes only when each term vanishes. Over $\mathbb{C}$ with the plain
transpose the statement is false. Take
\[
  A = \begin{pmatrix} 1 \\ i \end{pmatrix} \in \mathbb{C}^{2 \times 1},
  \qquad
  A^{\top}\!A = 1 \cdot 1 + i \cdot i = 1 + (-1) = 0 .
\]
Then $\ker(A^{\top}\!A) = \mathbb{C}$, while $A \neq 0$ gives
$\ker(A) = \{0\}$. Replacing the transpose by the conjugate transpose
restores the result, since
$A^{*}A = 1 \cdot 1 + (-i)\cdot i = 2$ is invertible. State the field
explicitly when writing the proof.
\end{remark}

\begin{corollary}
By the rank-nullity theorem applied to $A$ and to $A^{\top}\!A$, both of
which act on $\mathbb{R}^{n}$,
\[
  \rank(A^{\top}\!A) = n - \dim \ker(A^{\top}\!A)
                     = n - \dim \ker(A)
                     = \rank(A).
\]
Consequently $A^{\top}\!A$ is invertible if and only if $A$ has full
column rank, and the normal equations $A^{\top}\!A x = A^{\top} b$ then
have a unique solution.
\end{corollary}

\begin{remark}
This corollary underwrites least squares estimation and reappears in
interior point methods, where each iteration solves a system built from
$A^{\top} D A$ with $D$ a positive diagonal weight matrix.
\end{remark}

\end{document}
